\section{Composing sequences without joining their payloads}
\label{sec:sequence-piece-tree}

A slice can be described by two bounds, but a concatenation may contain pieces
from several sources. BioSeqDB stores such SQL expressions as expanded
PostgreSQL values containing an immutable balanced tree of pieces. An
\emph{expanded value} is a PostgreSQL value kept in a form suited to operations
in memory. It can later be converted to the flat form required for storage.

The tree has three kinds of node.
\begin{itemize}
  \item A \textbf{source leaf} retains access to a source and describes an
        interval in that source.
  \item A \textbf{literal leaf} owns the bytes of a literal and the information
        needed to interpret them.
  \item A \textbf{branch} joins two children. It records their total length,
        piece count, and tree height.
\end{itemize}
Reverse and complement flags belong to the edges. A reverse flag makes the
reader visit children in the opposite order and read each selected leaf
backward. A complement flag changes the ranks as they are read.

Figure~\ref{fig:piece-tree} uses the source from the slice example. It forms
\[
  Z=S[2,6)\ \Vert\ \operatorname{reverse}(S[8,12))\ \Vert\ \texttt{AC},
\]
where \(\Vert\) means concatenation. The three pieces produce
\texttt{GTAA}, \texttt{ATTG}, and \texttt{AC}. The result is
\texttt{GTAAATTGAC}, with length \(4+4+2=10\).

\begin{figure}[H]
\centering
\begin{tikzpicture}[x=1cm,y=1cm,font=\footnotesize,
  branch/.style={circle,draw=black!60,fill=black!4,
    minimum size=.68cm,inner sep=1pt},
  leaf/.style={draw=seqNineBlue,fill=seqNineBlue!10,
    minimum width=2.15cm,minimum height=.7cm,align=center}]
  \node[anchor=west,font=\bfseries] at (0,5.9)
    {(a) One source supplies two pieces};
  \foreach \i/\s in {0/A,1/C,2/G,3/T,4/A,5/A,6/C,7/C,8/G,9/T,10/T,11/A}
    \node[seq nine cell,fill=black!3] at (1.35+.72*\i,5.15) {\texttt{\s}};
  \draw[seqNineBlue,thick] (2.43,4.77) rectangle (5.31,5.53);
  \draw[seqNineOrange,thick] (6.75,4.77) rectangle (9.63,5.53);
  \node[anchor=east] at (.9,5.15) {$S$};
  \node[text=seqNineBlue] at (3.87,4.45) {$[2,6)$};
  \node[text=seqNineOrange] at (8.19,4.45) {$[8,12)$};
  \node[branch] (r) at (6.2,3.5) {10};
  \node[leaf] (a) at (3,2.15) {$S[2,6)$};
  \node[branch] (b) at (9,2.15) {6};
  \node[leaf,draw=seqNineOrange,fill=seqNineOrange!12] (c)
    at (7.25,.72) {$S[8,12)$};
  \node[leaf,draw=seqNineGreen,fill=seqNineGreen!12] (d)
    at (11,.72) {literal \texttt{AC}};
  \draw[seq nine flow] (r)--(a);
  \draw[seq nine flow] (r)--(b);
  \draw[seq nine flow,seqNineOrange] (b)--node[above left] {reverse} (c);
  \draw[seq nine flow] (b)--(d);
  \draw[dashed,seqNineBlue] (3.3,2.55)--(3.3,4.1);
  \draw[dashed,seqNineOrange,rounded corners=3pt]
    (c.south)--(7.25,.18)--(13.3,.18)--(13.3,4.1)--(8.19,4.1);
  \node[anchor=west,align=left] at (10.15,3.55)
    {circle: total length\\rectangle: one piece};
  \node[anchor=west,font=\bfseries] at (0,-.23)
    {(b) The reader emits the pieces in logical order};
  \foreach \i/\s in {0/G,1/T,2/A,3/A}
    \node[seq nine blue] at (2.05+.72*\i,-1.0) {\texttt{\s}};
  \foreach \i/\s in {0/A,1/T,2/T,3/G}
    \node[seq nine orange] at (5.33+.72*\i,-1.0) {\texttt{\s}};
  \foreach \i/\s in {0/A,1/C}
    \node[seq nine green] at (8.61+.72*\i,-1.0) {\texttt{\s}};
  \node[anchor=east] at (1.45,-1.0) {$Z$};
  \node[text=seqNineBlue] at (3.13,-1.6) {source, forward};
  \node[text=seqNineOrange] at (6.41,-1.6) {source, backward};
  \node[text=seqNineGreen] at (8.97,-1.6) {owned literal};
\end{tikzpicture}
\caption{A piece tree describes a ten-symbol result without putting the two
source intervals into a new sequence buffer. Dashed lines connect source
leaves to their intervals. The orange edge reverses one piece.}
\label{fig:piece-tree}
\end{figure}

The same tree can be shared by later expressions because existing nodes do
not change. Concatenation joins roots and restores balance. Slicing splits at
the two bounds and retains the middle part, sharing unaffected nodes.
Reversing a complete tree changes the edge flags. Two reversals cancel, as do
two complements when the alphabet's complement operation maps each symbol
back to itself after two applications.

This structure explains both the benefit and its limits.
\begin{itemize}
  \item A new expression can share existing payloads while allocating only
        the metadata needed to describe the change.
  \item Splits and joins follow paths in the tree. They are not universally
        constant-time operations.
  \item Reading \(m\) result symbols still requires producing those \(m\)
        ranks, together with the work of locating and reading their pieces.
  \item First expanding a flat value can copy its packed bytes to establish
        ownership. Binding a text literal also has to validate its symbols.
\end{itemize}
The point of the tree is to avoid repeatedly copying a long sequence merely
to describe a new order or interval.

\subsection{The alphabet is part of the value}

A plain SQL sequence literal can remain \emph{unbound}: it contains text but
has not yet adopted an alphabet. A bound value has a known alphabet and
packing width. The distinction prevents operations from silently changing a
symbol's meaning.
\begin{itemize}
  \item Slicing and reversal can operate on unbound text because they need
        only positions and order.
  \item Joining a bound sequence with a text literal validates the literal
        against that alphabet.
  \item Reverse-complement requires a bound alphabet with a complement rule
        for every symbol in that alphabet.
  \item Two different bound alphabets are not silently converted into one
        another.
\end{itemize}

Equality and hashing use logical content. For bound values this includes the
alphabet and the ordered ranks. Storage width, padding, and piece boundaries
do not change equality. An unbound text value has a different identity from a
bound sequence, even if both print the same letters. This is needed so that
SQL comparison and hash-based operations agree for all physical forms.
